% results.tex -- Results

We ran stochastic attention on all seven Pfam families and observed that the method consistently generated novel protein sequences while preserving family-level amino acid statistics, across a ten-fold range of family sizes and a seven-fold range of sequence lengths (Table~\ref{tab:results}, Fig.~\ref{fig:cross-family}).

We first examined whether the sampler converged to the correct Boltzmann distribution by comparing ULA and MALA outputs. MALA acceptance rates exceeded 99.6\% in every family tested (range: 99.6\%--99.9\%), confirming that the ULA discretization bias was negligible at $\alpha{=}0.01$. The two samplers produced near-identical metrics: for example, in the RRM family, SA retrieval achieved novelty $0.243$ and KL divergence $0.107$, while MALA yielded $0.249$ and $0.112$, respectively. This agreement held across all seven families, which justified the use of the computationally cheaper ULA in subsequent experiments.

We next assessed amino acid composition fidelity across families. In the generation regime ($\beta \approx 2\beta^*$), stochastic attention achieved KL divergences below $0.06$ in every family, with several families below $0.02$: WW achieved $\mathrm{KL}{=}0.008$, Kunitz $0.013$, and zinc finger $0.013$. These values were substantially lower than the bootstrap baseline ($\mathrm{KL}$ range: $0.007$--$0.133$), despite the fact that bootstrap samples are exact copies of stored patterns and thus trivially preserve composition at the individual-sequence level. The improvement arose because SA generation samples from the Boltzmann distribution over the entire family, producing a more representative ensemble than uniform resampling of a finite set. The convex combination baseline performed poorly, with KL divergences ranging from $0.06$ (Pkinase) to $2.35$ (zinc finger), which indicated that averaging across all stored patterns destroyed the peaked, position-specific residue distributions characteristic of conserved protein families (SI Appendix, Fig.~\ref{fig:aa-freq}).

We then examined whether the generated sequences were genuinely novel or merely replaying stored patterns. In the generation regime, novelty ranged from $0.42$ (PDZ) to $0.65$ (WW), indicating substantial departure from any single stored memory (Table~\ref{tab:results}). By contrast, bootstrap samples had zero novelty by construction, and Gaussian perturbation produced negligible novelty ($<0.02$ in all families), confirming that small additive noise is insufficient to escape the basin of the nearest stored pattern. The retrieval regime produced intermediate novelty ($0.17$--$0.35$), consistent with samples that explore locally around stored patterns without fully escaping their basins. Nearest sequence identity in amino acid space confirmed this picture: generation-regime sequences had $52$--$61$\% identity to the nearest stored sequence, comparable to the typical within-family identity, while retrieval-regime sequences remained at $53$--$85$\% identity depending on the family's intrinsic diversity.

We also characterized how the critical temperature varied across families. The empirical $\beta^*$ ranged from $3.2$ (PDZ, Pkinase) to $5.4$ (WW), consistently lower than the $\sqrt{d}$ prediction for random patterns ($5.8$--$13.6$). The ratio $\beta^*/\sqrt{d}$ was approximately $0.4$--$0.6$ across all families, which indicated that the structured similarity spectrum of real protein sequences, arising from conserved positions, correlated substitutions, and phylogenetic signal, shifts the retrieval-to-generation transition to lower temperatures. Families with higher mean column entropy ($\bar{H}_{\mathrm{col}}$) tended to have lower $\beta^*/\sqrt{d}$ ratios, consistent with the intuition that more diverse families require less thermal energy to escape stored patterns. The effective number of sequences $K_{\mathrm{eff}}$, computed by reweighting at 80\% identity, equaled $K$ in all but the WW family ($K_{\mathrm{eff}}{=}419$ vs.\ $K{=}420$), indicating minimal redundancy in the Pfam seed alignments.

We next turned to the scaling study, in which we subsampled the WW domain at $K \in \{20, 50, 100, 200, 400\}$ with five independent replicates per condition (Fig.~\ref{fig:scaling}). We observed that stochastic attention in the generation regime maintained low KL divergence across the entire range: $\mathrm{KL}{=}0.019 \pm 0.008$ at $K{=}20$, improving monotonically to $\mathrm{KL}{=}0.010 \pm 0.002$ at $K{=}400$. This graceful degradation contrasted sharply with the convex combination baseline, whose KL divergence exceeded $0.8$ at every family size and worsened with increasing $K$ ($0.81$ at $K{=}20$ to $1.50$ at $K{=}400$), confirming that averaging over more patterns compounds the smoothing artifact. Bootstrap and Gaussian perturbation maintained moderate KL values ($0.02$--$0.03$) but produced no novelty at any $K$, making them ineffective as generative methods.

The scaling study also revealed how PCA dimensionality and the critical temperature co-varied with family size. At $K{=}20$ the PCA reduction retained only $d{=}17$--$18$ components, producing a compact latent space in which $\beta^*{=}2.7$; at $K{=}400$ the richer alignment supported $d{=}182$--$183$ components and $\beta^*{=}5.5$. Despite this six-fold increase in effective dimensionality, SA generation quality remained stable, which suggested that the entropy inflection criterion successfully adapted the operating temperature to the geometry of the memory matrix at each scale. Novelty increased with $K$ (from $0.47$ at $K{=}20$ to $0.74$ at $K{=}400$), reflecting the richer landscape of a larger memory set, while nearest sequence identity decreased correspondingly (from $0.71$ to $0.56$), indicating that the sampler explored further from stored patterns when more patterns were available to collectively shape the energy landscape.

We finally asked whether the critical temperature could be predicted from alignment statistics alone, without running an explicit temperature sweep (Fig.~\ref{fig:beta-prediction}). We pooled 32 data points (the 7 Pfam families and 25 WW scaling replicates) and regressed $\beta^*$ against MSA-derived features. We found that a simple two-parameter model, $\beta^* \approx 1.52 + 0.28\sqrt{d}$ (intercept $1.52 \pm 0.08$, slope $0.282 \pm 0.007$; $\pm$ denotes bootstrap SE, $10{,}000$ resamples), explained $97\%$ of the variance ($R^2{=}0.969$, $\mathrm{RMSE}{=}0.16$), with $d$ the PCA dimension at $95\%$ retained variance. Adding further predictors (mean column entropy $\bar{H}_{\mathrm{col}}$, effective number of sequences $K_{\mathrm{eff}}$, or spectral concentration $\lambda_1/\mathrm{tr}(\mathbf{C})$) produced negligible improvement ($\Delta R^2 < 0.001$), which indicated that PCA dimensionality captured most of the relevant variation. The fitted slope of $0.28$ was roughly half the $\sqrt{d}$ coefficient predicted for random patterns, consistent with the empirical ratio $\beta^*/\sqrt{d} \approx 0.4$--$0.6$ observed across families. Importantly, the naive random-pattern prediction $\beta^* = \sqrt{d}$ yielded $\mathrm{RMSE}{=}5.3$, a $33\times$ larger error, confirming that the structured similarity spectrum of real protein families shifts the phase transition to substantially lower temperatures. This regression provides a practical shortcut: given a new family, one can estimate $\beta^*$ from PCA alone and set the generation temperature without a temperature sweep.

We compared stochastic attention against three established learned baselines: profile HMM emit, EvoDiff (MSA-conditioned diffusion), and the MSA Transformer (Gibbs MLM sampling) (Fig.~\ref{fig:baseline-comparison}). We used Wilcoxon rank-sum tests on per-chain estimates (30 chains $\times$ 5 samples) with Cohen's $d$ effect sizes. EvoDiff achieved the lowest KL divergence in most families ($0.003$--$0.011$), consistent with its pretraining on millions of UniRef50 MSAs, and profile HMM emit also produced low KL values ($0.006$--$0.030$) by directly capturing position-specific residue frequencies. SA generation was competitive with both ($0.008$--$0.060$) despite using only the family's seed alignment, typically $37$--$420$ sequences, with no external training data.

However, the three metrics together revealed a trade-off that no single learned baseline resolved. The critical metric was nearest sequence identity, where all three baselines differed significantly from SA generation ($p < 0.001$ in all families). HMM emit produced the highest novelty ($0.56$--$0.71$) but the lowest sequence identity to stored patterns ($11$--$42\%$; SA gen vs.\ HMM emit: $d = 3.6$--$24.6$ across families), indicating sequences so divergent from the family that they likely fall outside its functional envelope. The MSA Transformer showed a similar pattern: high novelty ($0.37$--$0.67$) but low sequence identity ($9$--$48\%$; $d = 2.1$--$26.8$). EvoDiff balanced KL and novelty more effectively but still produced significantly lower sequence identity than SA generation in every family tested ($p < 0.001$, $d = 1.6$--$5.5$), and its compute cost scaled prohibitively with sequence length: generating 150 sequences for the Pkinase family ($L{=}262$) required an estimated 18 hours on CPU, compared to seconds for SA. SA generation uniquely occupied the regime of low KL divergence, substantial novelty ($0.42$--$0.65$), and moderate sequence identity ($51$--$61\%$), generating sequences that departed meaningfully from stored patterns while remaining within the family's characteristic sequence space. Critically, SA achieved this balance without pretraining, backpropagation, or GPU resources, using only the information present in the seed alignment itself.

We assessed structural plausibility by predicting folded structures for 50 sequences per method per family using ESMFold (Fig.~\ref{fig:structure-validation}), comparing SA generation to stored (natural) sequences via Wilcoxon rank-sum tests with Cohen's $d$ effect sizes. SA-generated sequences achieved mean pLDDT scores statistically indistinguishable from stored natural sequences in five of seven families (Wilcoxon $p > 0.05$; RRM: $76.2$ vs.\ $76.9$, $p{=}0.45$; WW: $83.7$ vs.\ $82.9$, $p{=}0.08$; zf-C2H2: $83.4$ vs.\ $81.2$, $p{=}0.23$; PDZ: $77.6$ vs.\ $77.4$, $p{=}0.55$; Pkinase: $88.7$ vs.\ $89.7$, $p{=}0.11$). In Kunitz, SA generation achieved significantly \emph{higher} pLDDT than stored sequences ($90.8$ vs.\ $89.3$, $p{=}0.001$, $d{=}0.67$). Only SH3 showed lower pLDDT for SA generation ($70.5$ vs.\ $74.1$, $p{=}0.006$, $d{=}{-}0.56$), though the mean remained above the pLDDT${=}70$ confidence threshold. Across all families, $\geq 86\%$ of SA-generated sequences exceeded pLDDT${=}70$ (with $100\%$ in WW, Kunitz, and Pkinase).

TM-scores to family reference structures provided an even stronger result. In five of seven families, SA generation achieved significantly \emph{higher} TM-scores than stored sequences (Wilcoxon $p < 0.001$; RRM: $0.409$ vs.\ $0.380$, $d{=}1.62$; WW: $0.181$ vs.\ $0.174$, $d{=}1.14$; Kunitz: $0.847$ vs.\ $0.828$, $d{=}0.99$; PDZ: $0.638$ vs.\ $0.574$, $d{=}1.41$; Pkinase: $0.707$ vs.\ $0.679$, $d{=}1.56$). The remaining two families showed no significant difference (SH3: $p{=}0.56$; zf-C2H2: $p{=}0.11$). The large effect sizes ($d > 1.0$) in several families indicate that this is not merely a failure to reject the null: SA-generated sequences are structurally \emph{more similar} to the canonical family fold than a random sample of natural family members. This likely reflects the Boltzmann ensemble's tendency to weight the central tendency of the family distribution, producing sequences closer to the ``consensus fold'' than individual natural sequences, which may contain lineage-specific structural deviations. The three families with uniformly low TM-scores (RRM, WW, zf-C2H2) showed the same pattern for stored sequences, indicating a length mismatch between the generated domain and the multi-domain reference structure rather than a failure of the generative method. Taken together, these results demonstrate that stochastic attention generates sequences that are not merely statistically plausible strings of amino acids but encode three-dimensional folds consistent with, and in many cases more representative of, their target family, despite substantial sequence-level novelty ($0.42$--$0.65$).
